% ====================================================================
% 부록 A — 기호·부호 함정 검산표 (재수록 전용, 신규 자산 없음)
% ====================================================================
\section*{부록 A --- 기호$\cdot$부호 함정 검산표}
\addcontentsline{toc}{section}{부록 A --- 기호·부호 함정 검산표}

본문 여러 절에 흩어진 부호$\cdot$기호 함정을 한자리에 모은다. 각 항의 \emph{올바른 뜻}과 \emph{흔한 오류}는
해당 본문 절에서 명시했으며(근거 열), 이 표는 검수$\cdot$구현 시의 단일 조회점일 뿐 새 내용을 도입하지 않는다.
Chapter 1 부록의 부호 사슬 검산표와 같은 용도다.

\begin{longtable}{p{0.16\textwidth} p{0.30\textwidth} p{0.30\textwidth} c}
\toprule
항목 & 올바른 뜻 & 흔한 오류(피할 것) & 근거 \\
\midrule
\endfirsthead
\toprule
항목 & 올바른 뜻 & 흔한 오류(피할 것) & 근거 \\
\midrule
\endhead
$s$ vs $\sigma_d$ &
$s{=}+1$ 은 평형 관계~\eqref{eq:muV} 전용 고정 부호; $\sigma_d(\pm1)$ 는 방향(분극$\cdot$분기$\cdot$꼬리) 전용 &
$\sigma_d{=}-1$ 을 평형식에 기계 대입 $\to$ 점유/진행률 뒤바뀌는 \emph{여집합 오류}(종 대칭이라 잠복) &
\S\ref{ssec:logistic} \\
$\theta$ vs $\xi$ &
$\theta=$ Li 점유율(만충 $1$), $\xi=1-\theta=$ 탈리튬화 진행률(희박 $1$) --- 같은 분포의 두 이름
(\S\ref{ssec:logistic}) &
둘을 뒤섞으면 $\ln[\xi/(1-\xi)]$ 인자가 뒤집혀 발산 부호($\pm\infty$)를 반대로 읽음 &
\S\ref{ssec:dvdt} \\
$\bar x$ vs $x$ &
$\bar x\equiv1-x=$ 탈리튬화$\cdot$추출 전하 분율(음함수~\eqref{eq:implicit}$\cdot$그림~\ref{fig:blend} 가로축) &
Li 함량 $x$ 와 같은 배향으로 착각 $\to$ SOC 축 방향 반전 &
\S\ref{ssec:overlap} \\
$g$ 4종 &
$g_j(x){=}\xi(1-\xi)/w$ 봉우리 가중[1/V] &
$g_j(\xi)$(Ch1 조성 자유에너지)$\cdot$$g(\theta)$(BW)$\cdot$$g(E_F)$(상태밀도)와 혼동 &
\S\ref{ssec:overlap} \\
$F_\vib$ vs $F$ &
$F_\vib=$ 몰당 Helmholtz 자유에너지($\Delta F_\vib=RT\ln(1-e^{-\theta_E/T})$) &
Faraday 상수 $F$ 와 같은 기호로 오해 &
\S\ref{ssec:einstein-roundtrip} \\
$u_j$ vs $x$ &
$u_j\equiv\theta_{E,j}/T$ 무차원 진동 비 &
리튬 함량 $x$ 와 혼동 &
\S\ref{ssec:einstein-closed} \\
``방전($I>0$)'' &
Bernardi 셀-수지 관례 $=$ 흑연 하프셀 \emph{리튬화}; 부호는 라벨 아닌 전류 $I$ 로(충전 $I<0$) &
Ch1 ``방전 $=$ 탈리튬화'' 라벨과 \emph{같은 단어 반대 방향} 혼동 --- 최대 부호 위험 &
\S\ref{ssec:signlabel} \\
$\Delta S(x)$ vs $\Delta S_{a,j}$ &
반응(평형) 엔트로피 $\Delta S(x)$ 만 가역 발열에 들어감 &
활성화 엔트로피 $\Delta S_{a,j}$(비가역 동역학)를 섞음 &
\S\ref{ssec:threedist} \\
$\Delta S_\vib(T)$ vs $\partial S_\vib/\partial x$ &
$\Delta S_\vib(T){=}S_\vib(T){-}S_\vib(T_\mathrm{ref})$($x{,}\theta_E$ 고정, $T$ 편차) &
부분몰 $\partial S_\vib/\partial x$(고정 $T$, $x$ 편차)와 혼동 &
\S\ref{ssec:einstein-roundtrip} \\
$S_e$ 라벨 &
본 장 로컬 식은 \code{eq:Se-ch2}(자리당 Sommerfeld $S_e$, \S\ref{ssec:elec}) &
Chapter 1 §15 의 동명 식 \code{eq:Se}(별개 컴파일 객체)와 이름 혼동 &
\S\ref{ssec:elec} \\
\bottomrule
\end{longtable}

% 자산: (신규 자산 없음 — 본문 명시 CRITICAL 부호·기호 함정의 재수록 마스터표)
%   집약: [C-16] s vs σ_d · [C-21] θ/ξ · [C-72] x̄=1−x · [C-75] g 4종 · [C-57] F_vib/F ·
%   [C-50] u_j/x · [C-99] 방전(I>0) 라벨 · [C-71] 반응 vs 활성화 · [C-55] ΔS_vib(T)/∂S_vib/∂x
